\SourceLeaf{345}

\Needspace{8\baselineskip}\begin{center}
[III.]

\textbf{GEOMETRISCHE BEZIEHUNGEN MITHILFE KOMPLEXER ZAHLEN.}

I.
\end{center}

$M$, $M'$ zwei Punkte mit den komplexen Zahlen $m$, $m'$.

$P$ ein unbestimmter Punkt, für den die Azimute $PM'-PM$ einen vorgegebenen
Unterschied $= \varphi$ bilden.

Der geometrische Ort von $P$ ist ein Kreis mit

\[
\text{Mittelpunkt } \frac{m\,e^{i\varphi}-m'\,e^{-i\varphi}}{2i\sin\varphi},
\]

\[
\text{Radius } = \text{Betrag von } \frac{m'-m}{2\sin\varphi}.
\]

\Needspace{8\baselineskip}\begin{center}
II.
\end{center}

$M$, $M'$, $m$, $m'$ wie zuvor.

$P$ ein unbestimmter Punkt, dessen Abstände von $M$, $M'$ in einem gegebenen Verhältnis $k:k'$ stehen.

Der geometrische Ort von $P$ ist ein Kreis mit

\[
\text{Mittelpunkt } \frac{m.k'k'-m'.kk}{k'k'-kk} = c,
\]

\[
\text{Radius } = \text{Betrag von } \frac{kk'}{k'k'-kk}\cdot(m'-m).
\]



\Needspace{7\baselineskip}\begin{center}
\textbf{BEMERKUNGEN.}
\end{center}

Die drei vorstehenden Notizen stammen vermutlich aus der Zeit um 1840;
[I] steht auf einem einzelnen Blatt, [II] und [III] sind in einem Handbuch
eingetragen. Auf den Kreis bezieht sich auch eine bereits in Bd.\,IV.\,S.\,397 abgedruckte
»handschriftliche Notiz«. \hfill \textsc{Stäckel}.
